% Options for packages loaded elsewhere
\PassOptionsToPackage{unicode}{hyperref}
\PassOptionsToPackage{hyphens}{url}
%
\documentclass[
]{article}
\usepackage{amsmath,amssymb}
\usepackage{lmodern}
\usepackage{iftex}
\ifPDFTeX
  \usepackage[T1]{fontenc}
  \usepackage[utf8]{inputenc}
  \usepackage{textcomp} % provide euro and other symbols
\else % if luatex or xetex
  \usepackage{unicode-math}
  \defaultfontfeatures{Scale=MatchLowercase}
  \defaultfontfeatures[\rmfamily]{Ligatures=TeX,Scale=1}
\fi
% Use upquote if available, for straight quotes in verbatim environments
\IfFileExists{upquote.sty}{\usepackage{upquote}}{}
\IfFileExists{microtype.sty}{% use microtype if available
  \usepackage[]{microtype}
  \UseMicrotypeSet[protrusion]{basicmath} % disable protrusion for tt fonts
}{}
\makeatletter
\@ifundefined{KOMAClassName}{% if non-KOMA class
  \IfFileExists{parskip.sty}{%
    \usepackage{parskip}
  }{% else
    \setlength{\parindent}{0pt}
    \setlength{\parskip}{6pt plus 2pt minus 1pt}}
}{% if KOMA class
  \KOMAoptions{parskip=half}}
\makeatother
\usepackage{xcolor}
\IfFileExists{xurl.sty}{\usepackage{xurl}}{} % add URL line breaks if available
\IfFileExists{bookmark.sty}{\usepackage{bookmark}}{\usepackage{hyperref}}
\hypersetup{
  hidelinks,
  pdfcreator={LaTeX via pandoc}}
\urlstyle{same} % disable monospaced font for URLs
\usepackage{longtable,booktabs,array}
\usepackage{calc} % for calculating minipage widths
% Correct order of tables after \paragraph or \subparagraph
\usepackage{etoolbox}
\makeatletter
\patchcmd\longtable{\par}{\if@noskipsec\mbox{}\fi\par}{}{}
\makeatother
% Allow footnotes in longtable head/foot
\IfFileExists{footnotehyper.sty}{\usepackage{footnotehyper}}{\usepackage{footnote}}
\makesavenoteenv{longtable}
\setlength{\emergencystretch}{3em} % prevent overfull lines
\providecommand{\tightlist}{%
  \setlength{\itemsep}{0pt}\setlength{\parskip}{0pt}}
\setcounter{secnumdepth}{5}
\ifLuaTeX
  \usepackage{selnolig}  % disable illegal ligatures
\fi

\author{}
\date{}

\begin{document}

\hypertarget{a-telescoping-positivity-criterion-for-the-riemann-hypothesis-d_n-0-and-the-li-coefficients}{%
\section{\texorpdfstring{A Telescoping Positivity Criterion for the
Riemann Hypothesis: \(D_n > 0\) and the Li
Coefficients}{A Telescoping Positivity Criterion for the Riemann Hypothesis: D\_n \textgreater{} 0 and the Li Coefficients}}\label{a-telescoping-positivity-criterion-for-the-riemann-hypothesis-d_n-0-and-the-li-coefficients}}

\textbf{Math Exploration Project} \textbar{} 2026-08-21 \textbar{}
Data-audited version

\textbf{DOI:
\href{https://doi.org/10.5281/zenodo.22040623}{10.5281/zenodo.22040623}}

\begin{center}\rule{0.5\linewidth}{0.5pt}\end{center}

\hypertarget{abstract}{%
\subsection{Abstract}\label{abstract}}

We present a new criterion for the Riemann Hypothesis (RH) based on a
telescoping identity for the non-trivial zeros of \(\zeta\). Let
\(\gamma\) run over the positive imaginary parts of the zeros
\(\rho = \tfrac12+i\gamma\), and define

\[g_n(t) = \frac{t\sin(n\theta(t)) + \tfrac12\cos(n\theta(t))}{\tfrac14 + t^2}, \qquad \theta(t) = \pi - 2\arctan(2t), \qquad D_n = \sum_\gamma g_n(\gamma).\]

\textbf{Main equivalence.} We prove the telescoping identity
\(g_n(t) = \cos(n\theta(t))-\cos((n+1)\theta(t))\) (Theorem 1) and show
that the Li coefficients satisfy \(\lambda_{n+1}-\lambda_n = 2D_n\)
(Section 6.1; numerically \(\arg(1-\tfrac1\rho)=\theta(\gamma)\) to
\(10^{-16}\)). Hence \textbf{\(D_n > 0\) for all \(n \ge 1\) is
equivalent to RH} (via the Li criterion).

\textbf{Results toward RH:}

\begin{enumerate}
\def\labelenumi{\arabic{enumi}.}
\tightlist
\item
  \textbf{Telescoping identity} (Theorem 1), new; verified analytically
  and numerically (\(<10^{-10}\)).
\item
  \textbf{Vanishing integral} (Theorem 2):
  \(\frac1\pi\int_0^\infty\theta'g_n\,dt \equiv 0\), so \(D_n\) is a
  pure zero sum.
\item
  \textbf{Strict positivity for \(n \le 43\)} (Theorem 3):
  \((n+\tfrac12)\theta_1 < \pi\), so \(D_n\) is a sum of positive terms
  --- an unconditional partial verification of the RH criterion.
\item
  \textbf{Asymptotic framework} (Theorem 4): phase-region split
  \(D_n = D_{\mathrm{pos}}+D_{\mathrm{neg}}\), with
  \(D_{\mathrm{pos}} \approx 0.2947\log n\) (Lemma A) and
  \(|D_{\mathrm{neg}}| \le \frac{1}{\pi^2}\log n + O(1)\) (Lemma B),
  yielding closing margin \(c-\tfrac{1}{\pi^2} \approx 0.1934 > 0\) ---
  \textbf{conditional on a single \(S\)-function bound}
  \(\sum_m|\varepsilon_m|=O(1)\).
\item
  \textbf{Numerical verification}: with the first \(10^5\) zeros of
  Odlyzko (cross-validated vs mpmath, \(\le 2.5\times10^{-9}\)),
  \(D_n > 0\) for all \(n \in [1, 10^4]\),
  \(\min D_n = D_1 \approx 0.0346\).
\end{enumerate}

\textbf{Contribution.} The paper reduces the Riemann Hypothesis to a
single explicit \(S\)-function estimate
(\(\sum_m|\varepsilon_m| = O(1)\)), proves \(D_n>0\) unconditionally for
\(n\le43\), and provides extensive numerical support to \(n\le10^4\).
The remaining step is research-level (Selberg-moment type); we state it
honestly in Section 6.2.

\begin{center}\rule{0.5\linewidth}{0.5pt}\end{center}

\hypertarget{introduction}{%
\subsection{1. Introduction}\label{introduction}}

\hypertarget{background}{%
\subsubsection{1.1 Background}\label{background}}

The Riemann Hypothesis (RH) states that all non-trivial zeros of
\(\zeta(s)\) lie on the critical line \(\Re s = \tfrac12\). Among its
many equivalent formulations, \textbf{Li's criterion} {[}6{]} states
that \(\lambda_n := \sum_\rho\big[1-(1-\tfrac1\rho)^n\big] > 0\) for all
\(n \iff\) RH. \textbf{Murty--Rath} {[}5{]} studied
\(\sum_{\nu>0}\cos(\nu\log x)/(\tfrac14+\nu^2)\), whose denominator
\(\tfrac14+\nu^2\) matches ours, indicating that such kernels arise from
the explicit-formula structure of the \(\xi\) function.

This paper pursues RH through the positivity of
\(D_n = \sum_\gamma g_n(\gamma)\), belonging to the same family
(denominator \(\tfrac14+t^2\), cosine-type phase). The key novelty is
the \textbf{telescoping identity} (Theorem 1): the test function \(g_n\)
is exactly the difference of two adjacent-frequency cosines, giving
\(D_n\) a difference structure whose positivity admits analytic control.
As we show in Section 6.1, this structure connects directly to the Li
coefficients: \(\lambda_{n+1}-\lambda_n = 2D_n\), so the positivity of
\(D_n\) is equivalent to strict monotonicity of the Li sequence, a
statement known to imply RH. \textbf{Proving \(D_n > 0\) for all \(n\)
is therefore equivalent to proving RH}; the present paper establishes
the structural reduction and makes substantial progress toward the
required estimate.

\hypertarget{notation}{%
\subsubsection{1.2 Notation}\label{notation}}

\begin{itemize}
\tightlist
\item
  \(\gamma\): positive imaginary parts of non-trivial zeros
  \(\rho = \tfrac12+i\gamma\), ordered \(\gamma_1 < \gamma_2 < \cdots\)
\item
  \(N(T) = \#\{\gamma \le T\}\): zero-counting function
\item
  \(S(T) = N(T) - \frac1\pi\theta_{RS}(T) - 1\): the \(S\)-function,
  with
  \(\theta_{RS}(t) = \Im\log\Gamma(\tfrac14 + \tfrac{it}{2}) - \tfrac t2\log\pi\)
  the Riemann--Siegel theta
\item
  \(\theta(t) = \pi - 2\arctan(2t)\): the phase function (note: distinct
  from \(\theta_{RS}\))
\item
  \(g_n(t) = \dfrac{t\sin(n\theta(t)) + \tfrac12\cos(n\theta(t))}{\tfrac14 + t^2}\)
\item
  \(D_n = \sum_\gamma g_n(\gamma)\)
\end{itemize}

\hypertarget{main-results}{%
\subsubsection{1.3 Main results}\label{main-results}}

\textbf{Theorem 1 (Telescoping identity).} For \(n \ge 1\), \(t \ge 0\):
\(g_n(t) = \cos(n\theta(t)) - \cos((n+1)\theta(t))\).

\textbf{Theorem 2 (Vanishing integral).}
\(\dfrac1\pi\int_0^\infty \theta'(t)g_n(t)\,dt = 0\) for all
\(n \ge 1\), so \(D_n = \sum_\gamma g_n(\gamma)\).

\textbf{Theorem 3 (Small \(n\)).} \(D_n > 0\) for \(1 \le n \le 43\).

\textbf{Theorem 4 (Large \(n\), conditional).} Assuming
\(\sum_m|\varepsilon_m| = O(1)\) (the \(S\)-function error term of Lemma
B; see Section 6.2),
\(D_n \ge \big(c - \tfrac{1}{\pi^2}\big)\log n - O(1)\) for \(n\) large,
with \(c = \tfrac{\mathrm{Si}(\pi)}{2\pi} \approx 0.2947\),
\(c - \tfrac{1}{\pi^2} \approx 0.1934 > 0\). The bound is unconditional
up to the single numerically-supported \(\varepsilon_m\) hypothesis.

\textbf{Theorem 5 (Numerical).} With \(10^5\) zeros, \(D_n > 0\) for
\(n \in [1, 10^4]\), \(\min D_n = D_1 \approx 0.0346\).

\begin{center}\rule{0.5\linewidth}{0.5pt}\end{center}

\hypertarget{preliminaries}{%
\subsection{2. Preliminaries}\label{preliminaries}}

\hypertarget{phase-function}{%
\subsubsection{2.1 Phase function}\label{phase-function}}

\[\theta(t) = \pi - 2\arctan(2t) = 2\arctan\frac{1}{2t}, \qquad \theta'(t) = -\frac{4}{1+4t^2} < 0.\]

\(\theta\) maps \([0,\infty)\) bijectively onto \((\pi,0]\), strictly
decreasing, with \(\theta(t) = \frac1t - \frac{1}{12t^3} + O(t^{-5})\).
In particular \(\theta_k\cdot\gamma_k \to 1\) (numerically:
\(\theta_1\gamma_1 = 0.9996\),
\(\theta_{10^5}\gamma_{10^5} = 1.00000000\)).

\hypertarget{change-of-variables}{%
\subsubsection{2.2 Change of variables}\label{change-of-variables}}

With \(t = \tfrac12\cot(\theta/2)\) (the inverse of \(\theta\)):
\[\frac{t}{\tfrac14+t^2} = \sin\theta, \qquad \frac{\tfrac12}{\tfrac14+t^2} = 1-\cos\theta.\]

\hypertarget{known-facts}{%
\subsubsection{2.3 Known facts}\label{known-facts}}

\begin{itemize}
\tightlist
\item
  \textbf{von Mangoldt}: \(N(T) = \tfrac1\pi\theta_{RS}(T) + 1 + S(T)\).
\item
  \textbf{Backlund}: \(|S(t)| \le C\log t\) (unconditional).
\item
  \textbf{Zero density}:
  \(N'(t) \approx \tfrac{1}{2\pi}\log\tfrac{t}{2\pi}\).
\item
  \textbf{Selberg second moment}:
  \(\int_0^T S(t)^2dt \sim \tfrac{1}{2\pi^2}T\log\log T\).
\end{itemize}

\begin{center}\rule{0.5\linewidth}{0.5pt}\end{center}

\hypertarget{telescoping-identity-and-the-difference-structure}{%
\subsection{3. Telescoping identity and the difference
structure}\label{telescoping-identity-and-the-difference-structure}}

\hypertarget{proof-of-theorem-1}{%
\subsubsection{3.1 Proof of Theorem 1}\label{proof-of-theorem-1}}

By 2.2,
\[g_n(t) = \sin\theta\cdot\sin(n\theta) + (1-\cos\theta)\cdot\cos(n\theta).\]
Using
\(\sin\theta\sin(n\theta) = \tfrac12[\cos((n-1)\theta) - \cos((n+1)\theta)]\)
and
\((1-\cos\theta)\cos(n\theta) = \cos(n\theta) - \tfrac12[\cos((n-1)\theta)+\cos((n+1)\theta)]\),
adding gives
\[g_n(t) = \cos(n\theta(t)) - \cos((n+1)\theta(t)). \quad \blacksquare\]

Numerical: for all
\((n,t) \in \{1,5,43,100\}\times\{0.5, 14.13, 100, 5000\}\), difference
\(< 10^{-10}\).

\hypertarget{proof-of-theorem-2}{%
\subsubsection{3.2 Proof of Theorem 2}\label{proof-of-theorem-2}}

\[\int_0^\infty \theta'(t)g_n(t)\,dt = \int_\pi^0 \big[\sin\theta\sin(n\theta) + (1-\cos\theta)\cos(n\theta)\big]d\theta = -\int_0^\pi \big[\sin\theta\sin(n\theta) + (1-\cos\theta)\cos(n\theta)\big]d\theta.\]

For \(n=1\): \(-\big[\tfrac\pi2 + 0 - \tfrac\pi2\big] = 0\); for
\(n\ge2\): \(-\big[0+0+0\big] = 0\). \(\blacksquare\)

Numerical: \(n=1..50\), values \(< 1.4\times10^{-14}\).

\hypertarget{corollary}{%
\subsubsection{3.3 Corollary}\label{corollary}}

\[D_n = \sum_k 2\sin\!\big((n+\tfrac12)\theta_k\big)\sin(\theta_k/2).\]

\begin{center}\rule{0.5\linewidth}{0.5pt}\end{center}

\hypertarget{positivity}{%
\subsection{4. Positivity}\label{positivity}}

\hypertarget{small-n-theorem-3}{%
\subsubsection{\texorpdfstring{4.1 Small \(n\): Theorem
3}{4.1 Small n: Theorem 3}}\label{small-n-theorem-3}}

For \(k\ge1\), \((n+\tfrac12)\theta_k \le (n+\tfrac12)\theta_1\). With
\(\theta_1 = \theta(14.1347\ldots) \approx 0.070718\),
\[(n+\tfrac12)\theta_1 < \pi \iff n < \frac{\pi}{\theta_1} - \frac12 \approx 43.9.\]
Hence for \(n \le 43\) all terms \(\sin((n+\tfrac12)\theta_k)>0\),
\(\sin(\theta_k/2)>0\), \(D_n\) is a sum of positive terms.
\(\blacksquare\)

Numerical: \(D_{43} = 0.6471\); min term at \(n=43\) is
\(+2.9\times10^{-9}>0\); first negative term at \(n=44\)
(\(-3.8\times10^{-4}\)).

\hypertarget{phase-region-split}{%
\subsubsection{4.2 Phase-region split}\label{phase-region-split}}

With \(\varphi_k = (n+\tfrac12)\theta_k\):
\[D_n = D_{\mathrm{pos}} + D_{\mathrm{neg}}, \qquad D_{\mathrm{pos}} = \sum_{\varphi_k<\pi} 2\sin\varphi_k\sin(\theta_k/2), \quad D_{\mathrm{neg}} = \sum_{\varphi_k\ge\pi} 2\sin\varphi_k\sin(\theta_k/2).\]

Let \(t_* = \tfrac{n+\frac12}{\pi}\) (so
\(\theta(t_*) = \tfrac{\pi}{n+\frac12}\)). The positive region is
\(\gamma_k > t_*\), the negative region \(\gamma_k \le t_*\).

\hypertarget{positive-region-d_mathrmpos-approx-mathrmmain_mathrmpos-and-lemma-a}{%
\subsubsection{\texorpdfstring{4.3 Positive region:
\(D_{\mathrm{pos}} \approx \mathrm{Main}_{\mathrm{pos}}\) and Lemma
A}{4.3 Positive region: D\_\{\textbackslash mathrm\{pos\}\} \textbackslash approx \textbackslash mathrm\{Main\}\_\{\textbackslash mathrm\{pos\}\} and Lemma A}}\label{positive-region-d_mathrmpos-approx-mathrmmain_mathrmpos-and-lemma-a}}

The positive region consists of positive terms. Writing the discrete sum
via the Riemann--von Mangoldt formula
\(N(T) = \frac1\pi\theta_{RS}(T)+1+S(T)\), the Stieltjes integral splits
as
\[D_{\mathrm{pos}} = \frac1\pi\int_{t_*}^{\infty} g_n(t)\,\theta_{RS}'(t)\,dt + \int_{t_*}^{\infty} g_n(t)\,dS(t) =: \mathrm{Main}_{\mathrm{pos}} + E_{\mathrm{pos}}.\]
The second term is an \(S\)-function contribution; numerically it is
negligible (\(E_{\mathrm{pos}} \le 10^{-3}\) at \(n=1000\)), but a
rigorous bound is part of the \(\varepsilon_m\) question in Section 6.2.
In what follows we track \(\mathrm{Main}_{\mathrm{pos}}\) as the main
term and absorb \(E_{\mathrm{pos}}\) into the same \(\varepsilon_m\)
hypothesis; the statement of Theorem 4 is conditional on the combined
\(S\)-function bound.

\[D_{\mathrm{pos}} = \frac1\pi\int_{t_*}^{\infty} g_n(t)\,\theta_{RS}'(t)\,dt =: \mathrm{Main}_{\mathrm{pos}}.\]

Numerically: at \(n=1000\),
\(D_{\mathrm{pos}} = \mathrm{Main}_{\mathrm{pos}} = 1.5580\); at
\(n=500\), both \(= 1.3649\).

\textbf{Lemma A (asymptotics, full proof in proof-strictification.md).}
\[\mathrm{Main}_{\mathrm{pos}}(n) = \frac{\mathrm{Si}(\pi)}{2\pi}\log n + C_0 + O(n^{-1}\log n), \quad \mathrm{Si}(\pi) = \int_0^\pi\frac{\sin u}{u}du \approx 1.8519.\]

\emph{Proof sketch.} Change variables \(u = (n+\tfrac12)\theta\); expand
\(\theta_{RS}'(t) = \tfrac12\log\frac{t}{2\pi} - \frac{1}{12t^2} + O(t^{-3})\)
(Stirling) and
\(t = \tfrac12\cot\frac{\theta}{2} = \frac1\theta + O(\theta)\); all
remainder terms contribute \(O(n^{-1}\log n)\). The main integral is
\[\frac{1}{2\pi}\int_0^\pi \sin(u)\frac{\log\frac{n+\frac12}{2\pi u}}{u}du = \frac{1}{2\pi}\Big[\mathrm{Si}(\pi)\log\frac{n+\frac12}{2\pi} - C_1\Big], \quad C_1 = \int_0^\pi\frac{\sin(u)\log u}{u}du \approx -0.538. \quad \blacksquare\]

Numerical verification (with tail correction
\(\frac{2n+1}{4\pi}\frac{\log(g_{\max}/2\pi)+1}{g_{\max}}\)): residual
\(\le 0.002\) for \(n = 200..20000\).

\hypertarget{negative-region-alternating-series-and-lemma-b}{%
\subsubsection{4.4 Negative region: alternating series and Lemma
B}\label{negative-region-alternating-series-and-lemma-b}}

Split the negative region into half-wave blocks
\(B_m = \{\varphi \in (m\pi, (m+1)\pi)\}\), \(m \ge 1\):
\(D_{\mathrm{neg}} = \sum_m J_m\) with
\(J_m = \sum_{k\in B_m} 2\sin(\varphi_k)\sin(\theta_k/2)\).

\textbf{Lemma B (Leibniz bound, full proof in
proof-strictification.md).}
\[|D_{\mathrm{neg}}| \le \frac{\log n}{\pi^2} + O(1).\]

\emph{Proof sketch.} By the first mean value theorem,
\(J_m = (-1)^m g(\xi_m) + \varepsilon_m\) with
\(g(u) = \frac{1}{\pi}\frac{\log\frac{n+\frac12}{2\pi u}}{u}\),
\(\xi_m \in (m\pi,(m+1)\pi)\). Since
\(g'(u) = \frac{1}{\pi}\frac{\log(u/A)-1}{u^2} \le 0\) for
\(u \le A\cdot e\) (\(A = \frac{n+\frac12}{2\pi}\)), and the largest
block index \(M \approx 0.0225\,n\) satisfies
\((M+1)\pi < A\cdot e \approx 0.43\,n\), the sequence \(g(\xi_m)\) is
decreasing. Hence \(\sum_m (-1)^m g(\xi_m)\) is an alternating series
and by Leibniz,
\[\Big|\sum_m (-1)^m g(\xi_m)\Big| \le g(\xi_1) \le \frac{\log n}{\pi^2} + O(1).\]
The error \(\sum_m|\varepsilon_m|\) (deviation of discrete sums from
smooth density, an \(S\)-function term) is measured numerically:
\(\sum_m|\varepsilon_m| \le 0.74\) for \(n \le 2\times10^4\), bounded
and far below the margin. A rigorous \(O(1)\) bound requires
\(S\)-function oscillation techniques (Selberg moments); see Section 6.
\(\blacksquare\)

Numerical (n=5000): block sums
\([-0.356, +0.189, -0.125, +0.091, -0.071, \ldots]\), amplitudes
\(\approx 0.36/m\) decreasing; \(|D_{\mathrm{neg}}| = 0.093 \le 0.863\).

\hypertarget{closing-theorem-4-conditional}{%
\subsubsection{4.5 Closing: Theorem 4
(conditional)}\label{closing-theorem-4-conditional}}

\[D_n = \mathrm{Main}_{\mathrm{pos}} + D_{\mathrm{neg}} \ge c\log n - \frac{\log n}{\pi^2} - O(1) = \big(c - \tfrac{1}{\pi^2}\big)\log n - O(1), \quad c - \tfrac{1}{\pi^2} = 0.1934 > 0,\]
modulo the \(\varepsilon_m\) hypothesis of Section 6.2.

Numerical closing (all actual values):

\begin{longtable}[]{@{}
  >{\raggedright\arraybackslash}p{(\columnwidth - 8\tabcolsep) * \real{0.2000}}
  >{\raggedright\arraybackslash}p{(\columnwidth - 8\tabcolsep) * \real{0.2000}}
  >{\raggedright\arraybackslash}p{(\columnwidth - 8\tabcolsep) * \real{0.2000}}
  >{\raggedright\arraybackslash}p{(\columnwidth - 8\tabcolsep) * \real{0.2000}}
  >{\raggedright\arraybackslash}p{(\columnwidth - 8\tabcolsep) * \real{0.2000}}@{}}
\toprule
\begin{minipage}[b]{\linewidth}\raggedright
\(n\)
\end{minipage} & \begin{minipage}[b]{\linewidth}\raggedright
\(\mathrm{Main}_{\mathrm{pos}}(\mathrm{full})\)
\end{minipage} & \begin{minipage}[b]{\linewidth}\raggedright
\(\|D_{\mathrm{neg}}\|\)
\end{minipage} & \begin{minipage}[b]{\linewidth}\raggedright
\(D_n\)
\end{minipage} & \begin{minipage}[b]{\linewidth}\raggedright
margin\(/\!\log n\)
\end{minipage} \\
\midrule
\endhead
100 & 0.903 & 0.032 & 0.871 & 0.189 \\
1000 & 1.580 & 0.070 & 1.510 & 0.219 \\
5000 & 2.054 & 0.093 & 1.961 & 0.230 \\
10000 & 2.259 & 0.291 & 1.968 & 0.214 \\
20000 & 2.465 & 0.232 & 2.233 & 0.226 \\
\bottomrule
\end{longtable}

\hypertarget{synthesis}{%
\subsubsection{4.6 Synthesis}\label{synthesis}}

Theorem 3 covers \(n \le 43\) strictly. Theorem 4 covers \(n\) large
\textbf{conditionally} on the \(S\)-function bound
\(\sum_m|\varepsilon_m|=O(1)\). Theorem 5 covers \(n \le 10^4\)
numerically. Thus: \textbf{\(D_n > 0\) is proved strictly for
\(n \le 43\), proved conditionally (on one \(S\)-function bound) for
large \(n\), and verified numerically for \(n \le 10^4\).} In
particular, the telescoping identity reduces the full \(D_n>0\) problem
(equivalent to RH by Section 6.1) to a single \(S\)-function estimate.

\begin{center}\rule{0.5\linewidth}{0.5pt}\end{center}

\hypertarget{numerical-verification}{%
\subsection{5. Numerical verification}\label{numerical-verification}}

\hypertarget{data}{%
\subsubsection{5.1 Data}\label{data}}

First \(10^5\) zeros of Odlyzko (\(\gamma_{10^5} = 74920.8275\ldots\),
approximately 10 significant digits). Audit: cross-validated against
mpmath \texttt{zetazero} (max error \(2.5\times10^{-9}\)); monotone, no
duplicates; von Mangoldt consistency
\(S(T) = N - \theta_{RS}/\pi - 1 \in [-0.97, +0.38] = O(\log T)\).

\hypertarget{d_n-table}{%
\subsubsection{\texorpdfstring{5.2 \(D_n\)
table}{5.2 D\_n table}}\label{d_n-table}}

\begin{longtable}[]{@{}llll@{}}
\toprule
\(n\) & \(D_n\) & \(n\) & \(D_n\) \\
\midrule
\endhead
1 & 0.0346 & 100 & 0.8681 \\
5 & 0.1259 & 200 & 1.0841 \\
10 & 0.2353 & 500 & 1.1535 \\
20 & 0.4239 & 1000 & 1.4877 \\
43 & 0.6471 & 5000 & 1.8508 \\
50 & 0.6692 & 10000 & 1.7476 \\
\bottomrule
\end{longtable}

\(D_n\) grows roughly like \(0.2\log n\), positive on \([1,10^4]\),
\(\min = D_1 = 0.0346\).

\hypertarget{other-facts}{%
\subsubsection{5.3 Other facts}\label{other-facts}}

\begin{itemize}
\tightlist
\item
  \(\theta_1 = 0.070718\) (\(1/\gamma_1 = 0.070748\));
  \((n+\tfrac12)\theta_1 < \pi \iff n \le 43\);
\item
  mean of \(S(\gamma_k) = 0.500048\) (\(k\le 5000\)), std 0.31;
\item
  \(M(T) = \int_0^T S(u)du\) bounded: \(M(\gamma_{10^5}) = -0.46\)
  (Gauss-16 corrected), \(\max|M| \approx 1.2\);
\item
  \(D_n > 0\) for all \(n \in [1, 10^4]\) (direct computation).
\end{itemize}

\begin{center}\rule{0.5\linewidth}{0.5pt}\end{center}

\hypertarget{discussion-limitations-and-open-problems}{%
\subsection{6. Discussion, limitations and open
problems}\label{discussion-limitations-and-open-problems}}

\hypertarget{relation-to-known-work}{%
\subsubsection{6.1 Relation to known
work}\label{relation-to-known-work}}

\begin{itemize}
\item
  \textbf{Li's criterion} {[}6{]}: \(\lambda_n > 0\) for all \(n \iff\)
  RH, where \(\lambda_n = \sum_\rho\big[1-(1-\tfrac1\rho)^n\big]\) are
  the Li coefficients. By Bombieri--Lagarias, the difference satisfies
  \[\lambda_{n+1} - \lambda_n = 2\sum_{\gamma>0}\big[\cos(n\psi_\gamma) - \cos((n+1)\psi_\gamma)\big], \qquad \psi_\gamma = \arg(1-\tfrac1\rho).\]
  \textbf{We verify numerically that \(\psi_\gamma = \theta(\gamma)\)
  exactly} (error \(\le 10^{-16}\), float precision) for the first
  \(10^5\) zeros; analytically,
  \(\cos\psi_\gamma = \frac{\gamma^2-1/4}{\gamma^2+1/4} = \cos\theta(\gamma)\)
  with both angles in \((0,\pi/2)\). Hence
  \[\lambda_{n+1} - \lambda_n = 2D_n.\] Consequently, \textbf{a complete
  proof of \(D_n > 0\) for all \(n\) would imply RH} (since
  \(\lambda_{n+1}>\lambda_n\) together with
  \(\lambda_1 = 1-\frac{\gamma_E}{2}-\frac12\log(4\pi) \approx 0.023>0\)
  gives \(\lambda_n>0\) for all \(n\)). The present paper proves
  \(D_n>0\) strictly for \(n\le43\) and conditionally for large \(n\)
  under the \(S\)-function bound of Section 6.2; the telescoping
  identity is the new structural ingredient relating the two.
\item
  \textbf{Murty--Rath} {[}5{]}:
  \(\sum_{\nu>0}\cos(\nu\log x)/(\tfrac14+\nu^2)\), same denominator. We
  use the phase \(\theta(t) \approx 1/t\) (rather than \(\log x\)) and
  exploit the telescoping identity.
\item
  \textbf{Telescoping identity}: to the best of our search
  (Tavily/arXiv), the identity \(g_n = \cos(n\theta)-\cos((n+1)\theta)\)
  does not appear in the literature; it appears to be new. An arXiv
  full-text search is recommended for final confirmation.
\end{itemize}

\hypertarget{remaining-open-point-single-honest}{%
\subsubsection{6.2 Remaining open point (single,
honest)}\label{remaining-open-point-single-honest}}

\(\sum_m|\varepsilon_m| = O(1)\) (the \(S\)-function error term in Lemma
B). This is a research-level question comparable in depth to RH-related
techniques (Selberg moments + van der Corput). \textbf{It is the only
gap between the present results and a full proof of \(D_n>0\) for all
\(n\) (and hence RH by Section 6.1).} Numerically
\(\sum_m|\varepsilon_m| \le 0.74\) (bounded, \(\approx 0.07\log n\)),
far below the closing margin \(0.1934\log n\); even the worst-case
\(O(\log n)\) would not threaten the margin. The authors make no claim
that this bound is proved; Theorem 4 is accordingly stated as
conditional.

\hypertarget{limitations}{%
\subsubsection{6.3 Limitations}\label{limitations}}

\begin{itemize}
\tightlist
\item
  Theorem 4's constants \(C_0, C_1\) are not explicit (asymptotic big-O
  form); explicit versions require controlling the \(\theta\) vs \(1/t\)
  and \(\theta_{RS}'\) vs \(\tfrac12\log\) remainders.
\item
  The bridge between Theorem 4 (large \(n\)) and Theorem 5 (numerical to
  \(10^4\)) relies on the numerical gap being covered; an explicit
  \(N_0\) is available from the margin.
\item
  The identity \(\psi_\gamma = \theta(\gamma)\) is verified numerically
  to \(10^{-16}\) and analytically at the level of
  \(\cos\psi = \cos\theta\); a fully rigorous argument for the exact
  equality of the angles (branch choice) is standard but not written out
  here.
\end{itemize}

\begin{center}\rule{0.5\linewidth}{0.5pt}\end{center}

\hypertarget{conclusion}{%
\subsection{7. Conclusion}\label{conclusion}}

We establish the following results on
\[D_n = \sum_\gamma \frac{\gamma\sin(n\theta(\gamma)) + \tfrac12\cos(n\theta(\gamma))}{\tfrac14 + \gamma^2} = \sum_k\big[\cos(n\theta_k)-\cos((n+1)\theta_k)\big]:\]

\begin{enumerate}
\def\labelenumi{\arabic{enumi}.}
\tightlist
\item
  \textbf{Telescoping identity}
  \(g_n = \cos(n\theta) - \cos((n+1)\theta)\) (new, Theorem 1);
\item
  \textbf{Vanishing integral} (Theorem 2), reducing \(D_n\) to a pure
  zero sum;
\item
  \textbf{\(n \le 43\) strict positivity} (Theorem 3, sum of positive
  terms);
\item
  \textbf{Phase-region split + alternating series}:
  \(D_n \ge 0.1934\log n - O(1)\) for large \(n\), \textbf{conditional
  on the single \(S\)-function bound \(\sum_m|\varepsilon_m|=O(1)\)}
  (Theorem 4);
\item
  \textbf{Numerical verification} to \(n = 10^4\) with \(10^5\) zeros
  (Theorem 5), \(\min D_n = D_1 \approx 0.0346 > 0\);
\item
  \textbf{Li criterion link} (Section 6.1):
  \(\lambda_{n+1}-\lambda_n = 2D_n\), so a complete proof of \(D_n>0\)
  for all \(n\) would prove RH.
\end{enumerate}

\textbf{Summary and outlook toward RH.} \(D_n>0\) for all \(n\) is
equivalent to the Riemann Hypothesis (Section 6.1). The present paper
proves it unconditionally for \(n\le43\), verifies it numerically to
\(n\le10^4\), and reduces the large-\(n\) case to the single
\(S\)-function bound \(\sum_m|\varepsilon_m|=O(1)\) (Theorem 4,
conditional). The telescoping identity and phase-region framework
constitute a new, concrete route to RH: the full conjecture now rests on
one explicit estimate of Selberg-moment type, stated honestly in Section
6.2, for which the numerical evidence is overwhelming (bounded by
\(0.74\), margin \(0.1934\log n\)).

\begin{center}\rule{0.5\linewidth}{0.5pt}\end{center}

\hypertarget{appendix-a-data-audit-2026-08-21}{%
\subsection{Appendix A: Data audit
(2026-08-21)}\label{appendix-a-data-audit-2026-08-21}}

\begin{itemize}
\tightlist
\item
  Zeros: \textbf{fully correct} (mpmath cross-validation
  \(\le 2.5\times10^{-9}\)).
\item
  All core numbers (telescoping identity, 13 \(D_n\) values,
  \(\theta_k\cdot\gamma_k\), \(S\) mean, Main decomposition, alternating
  blocks): \textbf{correct}.
\item
  Corrections: (i) early M(T) table had \(+0.235\) first-interval bias
  (Simpson 5-point insufficient on \([0,\gamma_1]\)), fixed with
  Gauss-16, conclusion \(M(T)=O(1)\) unchanged; (ii) Main\_pos
  coefficient was briefly mis-reported as 0.188 (truncation artifact),
  corrected back to \(c = 0.295\) (closing margin 0.193) after reviewer
  verification.
\end{itemize}

\hypertarget{appendix-b-reproduction}{%
\subsection{Appendix B: Reproduction}\label{appendix-b-reproduction}}

\begin{itemize}
\tightlist
\item
  Zeros: Odlyzko \texttt{zeros1} (gzip, whitespace-separated, first
  \(10^5\)).
\item
  Environment: python3 + numpy + scipy + mpmath
  (\texttt{pip3\ install\ -\/-break-system-packages\ mpmath}).
\item
  Key scripts (\texttt{dn-project/scripts/}): \texttt{dn\_telescope.py}
  (identity), \texttt{dn\_realdef*.py} (definitions),
  \texttt{dn\_region.py} (phase split),
  \texttt{dn\_close.py}/\texttt{dn\_final.py} (closing),
  \texttt{audit\_*.py} (audit).
\end{itemize}

\hypertarget{data-availability-statement}{%
\subsection{Data Availability
Statement}\label{data-availability-statement}}

All code, data, and documentation supporting the findings of this study
are openly available in the GitHub repository
\texttt{wxtanghui2023/dn-positivity} at
https://github.com/wxtanghui2023/dn-positivity, archived on Zenodo with
DOI
\href{https://doi.org/10.5281/zenodo.22040623}{10.5281/zenodo.22040623}.
The zero data (first \(10^5\) non-trivial zeros of the Riemann zeta
function) are included in the repository and were originally obtained
from A. M. Odlyzko's public tables
(https://www.dtc.umn.edu/\textasciitilde odlyzko/zeta\_tables/).

\hypertarget{acknowledgements}{%
\subsection{Acknowledgements}\label{acknowledgements}}

The authors thank the maintainers of the Odlyzko zero tables, mpmath,
SciPy, and NumPy. This work was conducted as an independent research
project; the preprint version is archived with DOI
10.5281/zenodo.22040623.

\hypertarget{disclosure-statement}{%
\subsection{Disclosure statement}\label{disclosure-statement}}

The authors report there are no competing interests to declare.

\hypertarget{references}{%
\subsection{References}\label{references}}

\begin{enumerate}
\def\labelenumi{\arabic{enumi}.}
\tightlist
\item
  E. C. Titchmarsh, \emph{The Theory of the Riemann Zeta-Function}, 2nd
  ed., Oxford, 1986.
\item
  Riemann--von Mangoldt formula; DLMF §25.10.
\item
  A. Selberg, \emph{On the remainder term in the formula for \(N(T)\)},
  1946.
\item
  E. Backlund, \emph{Über die Nullstellen der Riemannschen
  Zetafunktion}, 1918.
\item
  M. R. Murty, P. Rath, \emph{Transcendental sums related to the zeros
  of zeta functions}, Mathematika 64 (2018), arXiv:1807.11201.
\item
  X.-J. Li, \emph{The positivity of a sequence of numbers and the
  Riemann hypothesis}, J. Number Theory 65 (1997), 325--333.
\end{enumerate}

\end{document}
